\documentclass[10pt,a4paper]{amsart}
\usepackage{amsmath,amssymb,amsthm,graphicx}
\usepackage{verbatim}
\newenvironment{reference}{\comment}{\endcomment}
%\newenvironment{reference}{}{}
\newenvironment{slogan}{\comment}{\endcomment}
\newenvironment{history}{\comment}{\endcomment}
% For commutative diagrams we use Xy-pic
\usepackage[all]{xy}

% We use 2cell for 2-commutative diagrams.
\xyoption{2cell}
% Theorem environments.
%
\theoremstyle{plain}
\newtheorem{theorem}[subsection]{Theorem}
\newtheorem{proposition}[subsection]{Proposition}
\newtheorem{lemma}[subsection]{Lemma}

\theoremstyle{definition}
\newtheorem{definition}[subsection]{Definition}
\newtheorem{example}[subsection]{Example}
\newtheorem{exercise}[subsection]{Exercise}
\newtheorem{situation}[subsection]{Situation}

\theoremstyle{remark}
\newtheorem{remark}[subsection]{Remark}
\newtheorem{remarks}[subsection]{Remarks}

\numberwithin{equation}{subsection}

% Macros
%
\def\lim{\mathop{\mathrm{lim}}\nolimits}
\def\colim{\mathop{\mathrm{colim}}\nolimits}
\def\Spec{\mathop{\mathrm{Spec}}}
\def\Hom{\mathop{\mathrm{Hom}}\nolimits}
\def\Ext{\mathop{\mathrm{Ext}}\nolimits}
\def\SheafHom{\mathop{\mathcal{H}\!\mathit{om}}\nolimits}
\def\SheafExt{\mathop{\mathcal{E}\!\mathit{xt}}\nolimits}
\def\Sch{\mathit{Sch}}
\def\Mor{\mathop{\mathrm{Mor}}\nolimits}
\def\Ob{\mathop{\mathrm{Ob}}\nolimits}
\def\Sh{\mathop{\mathit{Sh}}\nolimits}
\def\NL{\mathop{N\!L}\nolimits}
\def\CH{\mathop{\mathrm{CH}}\nolimits}
\def\proetale{{pro\text{-}\acute{e}tale}}
\def\etale{{\acute{e}tale}}
\def\QCoh{\mathit{QCoh}}
\def\Ker{\mathop{\mathrm{Ker}}}
\def\Im{\mathop{\mathrm{Im}}}
\def\Coker{\mathop{\mathrm{Coker}}}
\def\Coim{\mathop{\mathrm{Coim}}}

% Boxtimes
%
\DeclareMathSymbol{\boxtimes}{\mathbin}{AMSa}{"02}

%
% Macros for moduli stacks/spaces
%
\def\QCohstack{\mathcal{QC}\!\mathit{oh}}
\def\Cohstack{\mathcal{C}\!\mathit{oh}}
\def\Spacesstack{\mathcal{S}\!\mathit{paces}}
\def\Quotfunctor{\mathrm{Quot}}
\def\Hilbfunctor{\mathrm{Hilb}}
\def\Curvesstack{\mathcal{C}\!\mathit{urves}}
\def\Polarizedstack{\mathcal{P}\!\mathit{olarized}}
\def\Complexesstack{\mathcal{C}\!\mathit{omplexes}}
% \Pic is the operator that assigns to X its picard group, usage \Pic(X)
% \Picardstack_{X/B} denotes the Picard stack of X over B
% \Picardfunctor_{X/B} denotes the Picard functor of X over B
\def\Pic{\mathop{\mathrm{Pic}}\nolimits}
\def\Picardstack{\mathcal{P}\!\mathit{ic}}
\def\Picardfunctor{\mathrm{Pic}}
\def\Deformationcategory{\mathcal{D}\!\mathit{ef}}
\usepackage{cleveref}
\crefname{originallemma}{original lemma}{original lemmas}
\Crefname{originallemma}{Original lemma}{Original lemmas}
\makeatletter
\newlabel{lemma-flat-locally-finite-presentation}{{lemma-flat-locally-finite-presentation}{source}{}{sourcepointer-1}{}}
\newlabel{lemma-flat-locally-finite-presentation@cref}{{[originallemma][0][]lemma-flat-locally-finite-presentation}{[1][0][]source}{}{}{}}
\makeatother
\crefname{originallemma}{original lemma}{original lemmas}
\Crefname{originallemma}{Original lemma}{Original lemmas}
\makeatletter
\newlabel{lemma-construct-quasi-coherent-sheaves}{{lemma-construct-quasi-coherent-sheaves}{source}{}{sourcepointer-2}{}}
\newlabel{lemma-construct-quasi-coherent-sheaves@cref}{{[originallemma][0][]lemma-construct-quasi-coherent-sheaves}{[1][0][]source}{}{}{}}
\makeatother
\crefname{originalreference}{original reference}{original references}
\Crefname{originalreference}{Original reference}{Original references}
\makeatletter
\newlabel{more-algebra-section-determinants}{{more-algebra-section-determinants}{source}{}{sourcepointer-3}{}}
\newlabel{more-algebra-section-determinants@cref}{{[originalreference][0][]more-algebra-section-determinants}{[1][0][]source}{}{}{}}
\makeatother
\crefname{originalreference}{original reference}{original references}
\Crefname{originalreference}{Original reference}{Original references}
\makeatletter
\newlabel{more-algebra-remark-determinant-as-socle}{{more-algebra-remark-determinant-as-socle}{source}{}{sourcepointer-4}{}}
\newlabel{more-algebra-remark-determinant-as-socle@cref}{{[originalreference][0][]more-algebra-remark-determinant-as-socle}{[1][0][]source}{}{}{}}
\makeatother
\title{The determinant of a flat finitely presented module, defined as the annihilator inside its exterior algebra, is invertible}
\author{The Stacks Project Authors}
\date{stacks-project:modules.tex:local-snapshot-20261012; GNU FDL 1.2 or later}
\begin{document}
\maketitle
\noindent\textit{Editorial note.} All delivered mathematical text is the author's own, copied from the pinned Stacks Project chapter source sources/stacks-modules/pinned-original-source.tex; no statement or proof step was rewritten, reordered or removed.\par\medskip
\noindent\textit{Editorial note.} Preamble provenance: the Stacks Project keeps one single shared preamble for every chapter; this card takes it directly from the staged sources/stacks-modules/preamble.tex as three verbatim spans (lines 8-12, 14-18 and 154-217), so no derived file was written and nothing was added or redefined. The wrapper class and the base packages are supplied by the assembler.\par\medskip
\noindent\textit{Editorial note.} Bibliography provenance: neither the counted statement nor its proof contains a citation (this lane rescans the whole span before writing the card), so the card carries no bibliography and no entry was reconstructed.\par\medskip
\noindent\textit{Editorial note.} Selection provenance: this card was opened under the file-based control contract (CONTROL.md), with the mechanical pre-card gate run on its statement and proof interval before the card was written.\par\medskip
\noindent\textit{Allowed standard prerequisites (editorial).}\begin{itemize}
\item The local splitting of a flat module of finite presentation as a direct summand of a finite free module. Invoked at the start of the counted proof (source lines 4769-4771) as a pointer to a lemma of this chapter. Reduces the statement to modules which are direct summands of finite free modules.
\item The description of a quasi-coherent module associated to a module over the ring of global sections, and the algebraic description of the determinant of a finite projective module as the annihilator of the module inside its exterior algebra. Both are invoked inside the counted proof (source lines 4778-4787) as pointers to a lemma of this chapter and to results of another chapter. Transfers the algebraic statement to the sheaf level and identifies the two descriptions of the determinant.
\end{itemize}\medskip
\ifdefined\hypertarget\hypertarget{sourcepointer-1}{}\fi
\noindent\textit{Source reference.} Original lemma lemma-flat-locally-finite-presentation: sources/stacks-modules/pinned-original-source.tex, source lines 3041--3041. Navigation pointer only: the counted text refers to the Stacks item lemma-flat-locally-finite-presentation, which belongs to another chapter of the project and is not part of this card. The pointer carries the target's own label as its number because this chapter's staged snapshot does not contain the other chapter's tag table; its content is not reproduced here.\par\medskip
\ifdefined\hypertarget\hypertarget{sourcepointer-2}{}\fi
\noindent\textit{Source reference.} Original lemma lemma-construct-quasi-coherent-sheaves: sources/stacks-modules/pinned-original-source.tex, source lines 1425--1425. Navigation pointer only: the counted text refers to the Stacks item lemma-construct-quasi-coherent-sheaves, which belongs to another chapter of the project and is not part of this card. The pointer carries the target's own label as its number because this chapter's staged snapshot does not contain the other chapter's tag table; its content is not reproduced here.\par\medskip
\ifdefined\hypertarget\hypertarget{sourcepointer-3}{}\fi
\noindent\textit{Source reference.} Original reference more-algebra-section-determinants: sources/stacks-modules/pinned-original-source.tex, source lines 4785--4785. Navigation pointer only: the counted text refers to the Stacks item more-algebra-section-determinants, which belongs to another chapter of the project and is not part of this card. The pointer carries the target's own label as its number because this chapter's staged snapshot does not contain the other chapter's tag table; its content is not reproduced here.\par\medskip
\ifdefined\hypertarget\hypertarget{sourcepointer-4}{}\fi
\noindent\textit{Source reference.} Original reference more-algebra-remark-determinant-as-socle: sources/stacks-modules/pinned-original-source.tex, source lines 4786--4786. Navigation pointer only: the counted text refers to the Stacks item more-algebra-remark-determinant-as-socle, which belongs to another chapter of the project and is not part of this card. The pointer carries the target's own label as its number because this chapter's staged snapshot does not contain the other chapter's tag table; its content is not reproduced here.\par\medskip
\begin{lemma}
\label{lemma-determinant-as-socle}
Let $(X, \mathcal{O}_X)$ be a ringed space. Let $\mathcal{F}$
be a flat and finitely presented $\mathcal{O}_X$-module.
Denote
$$
\det(\mathcal{F}) \subset
\wedge^*_{\mathcal{O}_X}(\mathcal{F})
$$
the annihilator of $\mathcal{F} \subset \wedge^*_{\mathcal{O}_X}(\mathcal{F})$.
Then $\det(\mathcal{F})$ is an invertible $\mathcal{O}_X$-module.
\end{lemma}

\begin{proof}
To prove this we may work locally on $X$. Hence we may assume $\mathcal{F}$
is a direct summand of a finite free module, see
Lemma \ref{lemma-flat-locally-finite-presentation}. Say
$\mathcal{F} \oplus \mathcal{G} = \mathcal{O}_X^{\oplus n}$.
Set $R = \mathcal{O}_X(X)$. Then we see
$\mathcal{F}(X) \oplus \mathcal{G}(X) = R^{\oplus n}$
and correspondingly
$\mathcal{F}(U) \oplus \mathcal{G}(U) = \mathcal{O}_X(U)^{\oplus n}$
for all opens $U \subset X$.
We conclude that $\mathcal{F} = \mathcal{F}_M$ as in
Lemma \ref{lemma-construct-quasi-coherent-sheaves}
with $M = \mathcal{F}(X)$ a finite projective $R$-module.
In other words, we have $\mathcal{F}(U) = M \otimes_R \mathcal{O}_X(U)$.
This implies that
$\det(M) \otimes_R \mathcal{O}_X(U) = \det(\mathcal{F}(U))$
for all open $U \subset X$ with $\det$ as in
More on Algebra, Section \ref{more-algebra-section-determinants}. By
More on Algebra, Remark \ref{more-algebra-remark-determinant-as-socle}
we see that
$$
\det(M) \otimes_R \mathcal{O}_X(U) =
\det(\mathcal{F}(U)) \subset
\wedge^*_{\mathcal{O}_X(U)}(\mathcal{F}(U))
$$
is the annihilator of $\mathcal{F}(U)$. We conclude that
$\det(\mathcal{F})$ as defined in the statement of the lemma
is equal to $\mathcal{F}_{\det(M)}$. Some details omitted; one has
to be careful as annihilators cannot be defined as the sheafification
of taking annihilators on sections over opens. Thus $\det(\mathcal{F})$
is the pullback of an invertible module and we conclude.
\end{proof}

\end{document}
